\documentclass[11pt,a4paper]{article}
\usepackage[utf8]{inputenc}
\usepackage[T1]{fontenc}
\usepackage[catalan]{babel}
\usepackage{lmodern}
\usepackage[sfdefault,lining]{FiraSans}
\usepackage[a4paper,top=2.3cm,bottom=2.2cm,left=1.8cm,right=1.8cm,headheight=1.3cm,headsep=0.4cm,footskip=1cm]{geometry}
\usepackage{amsmath,amssymb}
\usepackage{icomma}
\usepackage{graphicx}
\usepackage[table]{xcolor}
\usepackage{tikz}
\usepackage{fontawesome5}
\usepackage{ccicons}
\usepackage{fancyhdr}
\usepackage{lastpage}
\usepackage{tasks}
\usepackage{needspace}
\usepackage[most]{tcolorbox}
\usepackage{array}
\usepackage[hidelinks,pdftitle={Solucions de la fitxa de la unitat 7: Estudi de funcions. 4t d'ESO Matemàtiques B},pdfsubject={Matemàtiques B 4t ESO},pdfkeywords={funcions, imatge, antiimatge, domini, recorregut, punts de tall, creixement, màxims i mínims, continuïtat, simetria, periodicitat, funció inversa}]{hyperref}

\definecolor{navy}{HTML}{1B2748}
\definecolor{accent}{HTML}{E8603A}
\definecolor{teal}{HTML}{1F9E90}
\definecolor{soft}{HTML}{EEF1F8}
\definecolor{muted}{HTML}{5D6685}

\setlength{\parindent}{0pt}
\setlength{\parskip}{3pt}
\settasks{label=\alph*), label-format=\bfseries\color{navy}, label-width=1.3em, item-indent=1.8em, column-sep=0.6em, after-item-skip=5pt, before-skip=3pt, after-skip=4pt}

% ---------- capçalera i peu ----------
\newcommand{\llicencia}{\href{https://creativecommons.org/licenses/by-nc-sa/4.0/deed.ca}{\ccbyncsa}\ \ Llicència Creative Commons Reconeixement-NoComercial-CompartirIgual 4.0 Internacional (CC BY-NC-SA 4.0)}
\pagestyle{fancy}
\fancyhf{}
\fancyhead[L]{\raisebox{-0.35\height}{\includegraphics[height=0.95cm]{logo-gva.png}}}
\fancyhead[C]{\small\color{muted}Solucions de la fitxa · Unitat 7: Estudi de funcions · 4t ESO B}
\fancyhead[R]{\raisebox{-0.35\height}{\includegraphics[height=0.95cm]{logo-vermellar.png}}}
\fancyfoot[L]{\scriptsize\color{muted}\llicencia}
\fancyfoot[R]{\small\color{muted}\thepage}
\renewcommand{\headrulewidth}{0.4pt}
\renewcommand{\headrule}{\color{muted!40}\hrule height \headrulewidth}
\fancypagestyle{primera}{\fancyhf{}\renewcommand{\headrulewidth}{0pt}
  \fancyfoot[L]{\scriptsize\color{muted}\llicencia}
  \fancyfoot[R]{\small\color{muted}\thepage}}

% ---------- nivells i exercicis ----------
\newcommand{\estrelles}[1]{{\color{accent}\ifcase#1\or\faStar\faStar[regular]\faStar[regular]\or\faStar\faStar\faStar[regular]\or\faStar\faStar\faStar\fi}}
\newcounter{exer}
\newcommand{\ex}[1]{\par\needspace{4\baselineskip}\medskip\refstepcounter{exer}\noindent{\color{navy}\bfseries\theexer.}\ \ {\small\estrelles{#1}}\ \ }

\newcommand{\seccio}[2]{\par\needspace{16\baselineskip}\bigskip
  \begin{tcolorbox}[colback=navy,colframe=navy,coltext=white,arc=2mm,boxrule=0pt,left=3mm,right=3mm,top=1.2mm,bottom=1.2mm]
  {\large\bfseries #1\quad #2}
  \end{tcolorbox}\vspace{-2pt}}

\newcommand{\exf}[1]{\par\medskip\refstepcounter{exer}\noindent{\color{navy}\bfseries\theexer.}\ \ {\small\estrelles{#1}}\ \ }
\newcolumntype{C}[1]{>{\centering\arraybackslash}p{#1}}
% quadre de recordatori (només a la fitxa d'exercicis)
\newcommand{\recorda}[1]{\par\smallskip\begin{tcolorbox}[colback=accent!7,colframe=accent!55,arc=2mm,boxrule=0.5pt,left=2.5mm,right=2.5mm,top=1mm,bottom=1mm]\small{\color{accent}\faLightbulb[regular]}\ \textbf{Recorda:}\ #1\end{tcolorbox}}
\usetikzlibrary{calc}
\usepackage{enumitem}
\definecolor{violet}{HTML}{7C3AED}
\definecolor{orange}{HTML}{EA8A1F}
\newcommand{\res}[1]{\boxed{#1}}\newcommand{\rest}[1]{\fbox{#1}}

\def\fmneed{11}% que el títol d'una secció no quede sol al final del full
% =====================================================================
% Capçalera FentMates de les fitxes: un color per a cada curs.
% Els colors són els de la pàgina inicial (LEVELS a assets/js/materials.js);
% el logo ve de _eines/logo/pdf/ (python3 _eines/logo/genera.py).
%
% Ús, al final del preàmbul de la unitat:
%   % =====================================================================
% Capçalera FentMates de les fitxes: un color per a cada curs.
% Els colors són els de la pàgina inicial (LEVELS a assets/js/materials.js);
% el logo ve de _eines/logo/pdf/ (python3 _eines/logo/genera.py).
%
% Ús, al final del preàmbul de la unitat:
%   % =====================================================================
% Capçalera FentMates de les fitxes: un color per a cada curs.
% Els colors són els de la pàgina inicial (LEVELS a assets/js/materials.js);
% el logo ve de _eines/logo/pdf/ (python3 _eines/logo/genera.py).
%
% Ús, al final del preàmbul de la unitat:
%   \input{../comu/fentmates.tex}
%   \fmcurs{1eso}
% i, al començament del document:
%   \capcalera{Matemàtiques}{1r ESO}{Fitxa d'exercicis · Unitat 2}{Divisibilitat}
% Opcional abans de \input: \def\fmneed{14} (espai mínim abans de cada secció).
% =====================================================================
\makeatletter
\def\fm@logodir{../../logo/pdf}
\providecommand{\fmneed}{10}
\def\fm@cursnom{}

% clau del curs, color i color del text damunt del color
\newcommand{\fm@curs}[3]{\@namedef{fm@color@#1}{#2}\@namedef{fm@text@#1}{#3}}
\fm@curs{1eso}{E6B84A}{navy}
\fm@curs{2eso}{3B82F6}{white}
\fm@curs{3eso}{FF7A50}{white}
\fm@curs{4eso-a}{A78BFA}{white}
\fm@curs{4eso-b}{6D28D9}{white}
\fm@curs{1batx-mi}{9B1C1C}{white}
\fm@curs{1batx-ccss}{EF6B6B}{white}
\fm@curs{2batx-mii}{1E6B47}{white}
\fm@curs{2batx-ccss}{2BB3A3}{white}

\newcommand{\fmcurs}[1]{%
  \@ifundefined{fm@color@#1}{\PackageError{fentmates}{Curs desconegut: #1}{}}{}%
  \edef\fm@tmp{\noexpand\definecolor{curs}{HTML}{\@nameuse{fm@color@#1}}%
               \noexpand\colorlet{curstext}{\@nameuse{fm@text@#1}}}\fm@tmp
  \colorlet{cursfosc}{curs!72!black}%
  \def\fm@logo{\fm@logodir/vertical-#1.pdf}%
  \def\fm@marca{\fm@logodir/marca-#1.pdf}}

% ---------- capçalera de la primera pàgina ----------
% #1 matèria (petit), #2 curs, #3 tipus de document i unitat, #4 títol
\newcommand{\capcalera}[4]{%
  \gdef\fm@cursnom{#2}%
  \par\noindent
  \begin{tcolorbox}[colback=curs,colframe=curs,boxrule=0pt,arc=2.2mm,width=3.5cm,height=2.15cm,
      valign=center,left=3mm,right=2mm,top=1mm,bottom=1mm,nobeforeafter,box align=center]
    \color{curstext}{\footnotesize #1}\par\vspace{1pt}{\fontsize{19}{21}\selectfont\bfseries #2\par}%
  \end{tcolorbox}\hspace{5mm}%
  \begin{minipage}[c]{\dimexpr\linewidth-3.5cm-5mm-2.9cm\relax}
    {\color{cursfosc}\small\bfseries\MakeUppercase{#3}}\par\vspace{2pt}
    {\color{navy}\fontsize{20}{23}\selectfont\bfseries #4\par}
  \end{minipage}\hfill
  \begin{minipage}[c]{2.6cm}\raggedleft\includegraphics[width=2.6cm]{\fm@logo}\end{minipage}\par
  \vspace{2pt}}

% ---------- barres de secció del color del curs ----------
\renewcommand{\seccio}[2]{\par\needspace{\fmneed\baselineskip}\bigskip
  \begin{tcolorbox}[colback=curs,colframe=curs,coltext=curstext,arc=2mm,boxrule=0pt,left=3mm,right=3mm,top=1.2mm,bottom=1.2mm]
  {\large\bfseries #1\quad #2}
  \end{tcolorbox}\vspace{-2pt}}

% ---------- estrelles de dificultat del color del curs (a Batxillerat no se'n posen) ----------
\@ifundefined{estrelles}{}{%
  \renewcommand{\estrelles}[1]{{\color{curs}\ifcase#1\or\faStar\faStar[regular]\faStar[regular]\or\faStar\faStar\faStar[regular]\or\faStar\faStar\faStar\fi}}}

% ---------- capçalera i peu de les pàgines ----------
% primera pàgina: logos institucionals al peu
\fancypagestyle{primera}{\fancyhf{}\renewcommand{\headrulewidth}{0pt}
  \fancyfoot[L]{\raisebox{-0.32\height}{\includegraphics[height=0.78cm]{logo-gva.png}}\hspace{5mm}%
    \raisebox{-0.32\height}{\includegraphics[height=0.78cm]{logo-vermellar.png}}}
  \fancyfoot[R]{\scriptsize\color{muted}\href{https://creativecommons.org/licenses/by-nc-sa/4.0/deed.ca}{\ccbyncsa}\ \ CC BY-NC-SA 4.0 · fentmates.org\qquad\small\thepage}}
% la resta: el símbol a l'esquerra i el curs a la dreta, amb el títol al centre com abans
\fancyhead[L]{\raisebox{-0.3\height}{\includegraphics[height=0.62cm]{\fm@marca}}}
\fancyhead[R]{\small\bfseries\color{cursfosc}\fm@cursnom}
\renewcommand{\headrule}{\color{curs!45}\hrule height \headrulewidth}
\makeatother

%   \fmcurs{1eso}
% i, al començament del document:
%   \capcalera{Matemàtiques}{1r ESO}{Fitxa d'exercicis · Unitat 2}{Divisibilitat}
% Opcional abans de \input: \def\fmneed{14} (espai mínim abans de cada secció).
% =====================================================================
\makeatletter
\def\fm@logodir{../../logo/pdf}
\providecommand{\fmneed}{10}
\def\fm@cursnom{}

% clau del curs, color i color del text damunt del color
\newcommand{\fm@curs}[3]{\@namedef{fm@color@#1}{#2}\@namedef{fm@text@#1}{#3}}
\fm@curs{1eso}{E6B84A}{navy}
\fm@curs{2eso}{3B82F6}{white}
\fm@curs{3eso}{FF7A50}{white}
\fm@curs{4eso-a}{A78BFA}{white}
\fm@curs{4eso-b}{6D28D9}{white}
\fm@curs{1batx-mi}{9B1C1C}{white}
\fm@curs{1batx-ccss}{EF6B6B}{white}
\fm@curs{2batx-mii}{1E6B47}{white}
\fm@curs{2batx-ccss}{2BB3A3}{white}

\newcommand{\fmcurs}[1]{%
  \@ifundefined{fm@color@#1}{\PackageError{fentmates}{Curs desconegut: #1}{}}{}%
  \edef\fm@tmp{\noexpand\definecolor{curs}{HTML}{\@nameuse{fm@color@#1}}%
               \noexpand\colorlet{curstext}{\@nameuse{fm@text@#1}}}\fm@tmp
  \colorlet{cursfosc}{curs!72!black}%
  \def\fm@logo{\fm@logodir/vertical-#1.pdf}%
  \def\fm@marca{\fm@logodir/marca-#1.pdf}}

% ---------- capçalera de la primera pàgina ----------
% #1 matèria (petit), #2 curs, #3 tipus de document i unitat, #4 títol
\newcommand{\capcalera}[4]{%
  \gdef\fm@cursnom{#2}%
  \par\noindent
  \begin{tcolorbox}[colback=curs,colframe=curs,boxrule=0pt,arc=2.2mm,width=3.5cm,height=2.15cm,
      valign=center,left=3mm,right=2mm,top=1mm,bottom=1mm,nobeforeafter,box align=center]
    \color{curstext}{\footnotesize #1}\par\vspace{1pt}{\fontsize{19}{21}\selectfont\bfseries #2\par}%
  \end{tcolorbox}\hspace{5mm}%
  \begin{minipage}[c]{\dimexpr\linewidth-3.5cm-5mm-2.9cm\relax}
    {\color{cursfosc}\small\bfseries\MakeUppercase{#3}}\par\vspace{2pt}
    {\color{navy}\fontsize{20}{23}\selectfont\bfseries #4\par}
  \end{minipage}\hfill
  \begin{minipage}[c]{2.6cm}\raggedleft\includegraphics[width=2.6cm]{\fm@logo}\end{minipage}\par
  \vspace{2pt}}

% ---------- barres de secció del color del curs ----------
\renewcommand{\seccio}[2]{\par\needspace{\fmneed\baselineskip}\bigskip
  \begin{tcolorbox}[colback=curs,colframe=curs,coltext=curstext,arc=2mm,boxrule=0pt,left=3mm,right=3mm,top=1.2mm,bottom=1.2mm]
  {\large\bfseries #1\quad #2}
  \end{tcolorbox}\vspace{-2pt}}

% ---------- estrelles de dificultat del color del curs (a Batxillerat no se'n posen) ----------
\@ifundefined{estrelles}{}{%
  \renewcommand{\estrelles}[1]{{\color{curs}\ifcase#1\or\faStar\faStar[regular]\faStar[regular]\or\faStar\faStar\faStar[regular]\or\faStar\faStar\faStar\fi}}}

% ---------- capçalera i peu de les pàgines ----------
% primera pàgina: logos institucionals al peu
\fancypagestyle{primera}{\fancyhf{}\renewcommand{\headrulewidth}{0pt}
  \fancyfoot[L]{\raisebox{-0.32\height}{\includegraphics[height=0.78cm]{logo-gva.png}}\hspace{5mm}%
    \raisebox{-0.32\height}{\includegraphics[height=0.78cm]{logo-vermellar.png}}}
  \fancyfoot[R]{\scriptsize\color{muted}\href{https://creativecommons.org/licenses/by-nc-sa/4.0/deed.ca}{\ccbyncsa}\ \ CC BY-NC-SA 4.0 · fentmates.org\qquad\small\thepage}}
% la resta: el símbol a l'esquerra i el curs a la dreta, amb el títol al centre com abans
\fancyhead[L]{\raisebox{-0.3\height}{\includegraphics[height=0.62cm]{\fm@marca}}}
\fancyhead[R]{\small\bfseries\color{cursfosc}\fm@cursnom}
\renewcommand{\headrule}{\color{curs!45}\hrule height \headrulewidth}
\makeatother

%   \fmcurs{1eso}
% i, al començament del document:
%   \capcalera{Matemàtiques}{1r ESO}{Fitxa d'exercicis · Unitat 2}{Divisibilitat}
% Opcional abans de \input: \def\fmneed{14} (espai mínim abans de cada secció).
% =====================================================================
\makeatletter
\def\fm@logodir{../../logo/pdf}
\providecommand{\fmneed}{10}
\def\fm@cursnom{}

% clau del curs, color i color del text damunt del color
\newcommand{\fm@curs}[3]{\@namedef{fm@color@#1}{#2}\@namedef{fm@text@#1}{#3}}
\fm@curs{1eso}{E6B84A}{navy}
\fm@curs{2eso}{3B82F6}{white}
\fm@curs{3eso}{FF7A50}{white}
\fm@curs{4eso-a}{A78BFA}{white}
\fm@curs{4eso-b}{6D28D9}{white}
\fm@curs{1batx-mi}{9B1C1C}{white}
\fm@curs{1batx-ccss}{EF6B6B}{white}
\fm@curs{2batx-mii}{1E6B47}{white}
\fm@curs{2batx-ccss}{2BB3A3}{white}

\newcommand{\fmcurs}[1]{%
  \@ifundefined{fm@color@#1}{\PackageError{fentmates}{Curs desconegut: #1}{}}{}%
  \edef\fm@tmp{\noexpand\definecolor{curs}{HTML}{\@nameuse{fm@color@#1}}%
               \noexpand\colorlet{curstext}{\@nameuse{fm@text@#1}}}\fm@tmp
  \colorlet{cursfosc}{curs!72!black}%
  \def\fm@logo{\fm@logodir/vertical-#1.pdf}%
  \def\fm@marca{\fm@logodir/marca-#1.pdf}}

% ---------- capçalera de la primera pàgina ----------
% #1 matèria (petit), #2 curs, #3 tipus de document i unitat, #4 títol
\newcommand{\capcalera}[4]{%
  \gdef\fm@cursnom{#2}%
  \par\noindent
  \begin{tcolorbox}[colback=curs,colframe=curs,boxrule=0pt,arc=2.2mm,width=3.5cm,height=2.15cm,
      valign=center,left=3mm,right=2mm,top=1mm,bottom=1mm,nobeforeafter,box align=center]
    \color{curstext}{\footnotesize #1}\par\vspace{1pt}{\fontsize{19}{21}\selectfont\bfseries #2\par}%
  \end{tcolorbox}\hspace{5mm}%
  \begin{minipage}[c]{\dimexpr\linewidth-3.5cm-5mm-2.9cm\relax}
    {\color{cursfosc}\small\bfseries\MakeUppercase{#3}}\par\vspace{2pt}
    {\color{navy}\fontsize{20}{23}\selectfont\bfseries #4\par}
  \end{minipage}\hfill
  \begin{minipage}[c]{2.6cm}\raggedleft\includegraphics[width=2.6cm]{\fm@logo}\end{minipage}\par
  \vspace{2pt}}

% ---------- barres de secció del color del curs ----------
\renewcommand{\seccio}[2]{\par\needspace{\fmneed\baselineskip}\bigskip
  \begin{tcolorbox}[colback=curs,colframe=curs,coltext=curstext,arc=2mm,boxrule=0pt,left=3mm,right=3mm,top=1.2mm,bottom=1.2mm]
  {\large\bfseries #1\quad #2}
  \end{tcolorbox}\vspace{-2pt}}

% ---------- estrelles de dificultat del color del curs (a Batxillerat no se'n posen) ----------
\@ifundefined{estrelles}{}{%
  \renewcommand{\estrelles}[1]{{\color{curs}\ifcase#1\or\faStar\faStar[regular]\faStar[regular]\or\faStar\faStar\faStar[regular]\or\faStar\faStar\faStar\fi}}}

% ---------- capçalera i peu de les pàgines ----------
% primera pàgina: logos institucionals al peu
\fancypagestyle{primera}{\fancyhf{}\renewcommand{\headrulewidth}{0pt}
  \fancyfoot[L]{\raisebox{-0.32\height}{\includegraphics[height=0.78cm]{logo-gva.png}}\hspace{5mm}%
    \raisebox{-0.32\height}{\includegraphics[height=0.78cm]{logo-vermellar.png}}}
  \fancyfoot[R]{\scriptsize\color{muted}\href{https://creativecommons.org/licenses/by-nc-sa/4.0/deed.ca}{\ccbyncsa}\ \ CC BY-NC-SA 4.0 · fentmates.org\qquad\small\thepage}}
% la resta: el símbol a l'esquerra i el curs a la dreta, amb el títol al centre com abans
\fancyhead[L]{\raisebox{-0.3\height}{\includegraphics[height=0.62cm]{\fm@marca}}}
\fancyhead[R]{\small\bfseries\color{cursfosc}\fm@cursnom}
\renewcommand{\headrule}{\color{curs!45}\hrule height \headrulewidth}
\makeatother

\fmcurs{4eso-b}

\begin{document}
\thispagestyle{primera}

\capcalera{Matemàtiques B}{4t ESO}{Solucions de la fitxa · Unitat 7}{\raggedright Estudi de funcions}
\vspace{8pt}

\begin{tcolorbox}[colback=soft,colframe=soft,arc=2mm,left=3mm,right=3mm,top=1.5mm,bottom=1.5mm]
\small Intenta resoldre cada exercici abans de mirar-ne la solució. En els exercicis de «dibuixa una funció…» hi ha moltes respostes correctes: comprova que la teua complisca totes les condicions. Pots comprovar dominis, talls, simetries i inverses amb les eines interactives dels apunts de la unitat.
\end{tcolorbox}

% =====================================================
\seccio{1}{Concepte de funció. Imatges i antiimatges}

\exf{1} \begin{tasks}(4)
  \task $\res{\text{Sí}}$
  \task $\res{\text{No}}$
  \task $\res{\text{Sí}}$
  \task $\res{\text{No}}$
\end{tasks}Les gràfiques b) i d) tenen valors de $x$ amb dues imatges: hi ha rectes verticals que les tallen dues vegades.

\ex{1} \begin{tasks}(1)
  \task $\res{\text{Sí}}$
  \task $\res{\text{No}}$: el 4 té dues arrels, $2$ i $-2$.
  \task $\res{\text{Sí}}$
  \task $\res{\text{No}}$: un alumne pot tindre dos o més germans (o cap).
\end{tasks}

\ex{1} \begin{tasks}(1)
  \task $f(-2)=4+4-3=\res{5}$; \ $f(0)=\res{-3}$; \ $f\left(\frac12\right)=\frac14-1-3=\res{-\frac{15}{4}}$
  \task $x^2-2x-3=0\Rightarrow\res{x=-1,\ x=3}$; \ $x^2-2x+1=0\Rightarrow\res{x=1}$; \ $x^2-2x+3=0$, $\Delta=-8<0$: $-6$ \rest{no té antiimatges}.
\end{tasks}

\ex{1} \begin{center}\renewcommand{\arraystretch}{1.25}\begin{tabular}{|c|C{1cm}|C{1cm}|C{1cm}|C{1cm}|C{1cm}|}\hline $x$ & $-2$ & $-1$ & $0$ & $1$ & $2$\\\hline $f(x)$ & $\res{-8}$ & $\res{-5}$ & $\res{-2}$ & $\res{1}$ & $\res{4}$\\\hline\end{tabular}\end{center} $3x-2=10\Rightarrow\res{x=4}$

\ex{2} \begin{tasks}(4)
  \task $\res{f(x)=3x+5}$
  \task $\res{A(x)=x^2}$
  \task $\res{P(x)=1{,}2x}$
  \task $\res{f(x)=\dfrac1x}$
\end{tasks}

\ex{2} L'altre costat mesura $20-x$: $\res{A(x)=x(20-x)=20x-x^2}$. Cal $0<x<20$: $D=\res{(0,\,20)}$. \ $A(5)=\res{75\text{ m}^2}$, $A(10)=\res{100\text{ m}^2}$.


% =====================================================
\seccio{2}{Domini i recorregut}

\exf{1} \begin{tasks}(2)
  \task $D=\res{[-1,\,3]}$, $R=\res{[-2,\,2]}$
  \task $D=\res{\mathbb{R}-\{-1\}}$, $R=\res{\mathbb{R}-\{2\}}$
  \task $D=\res{[-2,\,+\infty)}$, $R=\res{[-1,\,+\infty)}$
  \task $D=\res{[-4,\,2]}$, $R=\res{[-1,\,2)\cup\{3\}}$
\end{tasks}

\ex{1} \begin{tasks}(3)
  \task $\res{\mathbb{R}}$
  \task $\res{\mathbb{R}-\{5\}}$
  \task $\res{\mathbb{R}-\{-3,\,3\}}$
  \task $\res{\mathbb{R}}$
  \task $\res{\mathbb{R}-\{-2,\,0,\,2\}}$
  \task $\res{\mathbb{R}-\{2,\,3\}}$
\end{tasks}

\ex{2} \begin{tasks}(3)
  \task $\res{[3,\,+\infty)}$
  \task $\res{[-3,\,3]}$
  \task $\res{(-\infty,\,1]\cup[3,\,+\infty)}$
  \task $\res{\mathbb{R}}$
  \task $\res{(2,\,+\infty)}$
  \task $\res{[-3,\,1)\cup(1,\,+\infty)}$
\end{tasks}

\ex{1} $f(2)=\sqrt{-1}$ no existeix: $\res{2\notin D(f)}$. \ $f(3)=\sqrt4=2$: $\res{3\in D(f)}$. \ $x^2-5\geq0\Rightarrow D(f)=\res{(-\infty,\,-\sqrt5\,]\cup[\,\,\sqrt5,\,+\infty)}$

\ex{2} \begin{tasks}(2)
  \task $x^2\geq0\Rightarrow R=\res{[2,\,+\infty)}$
  \task $\sqrt{x-1}\geq0\Rightarrow R=\res{[3,\,+\infty)}$
  \task $\frac2x$ no val mai 0: $R=\res{\mathbb{R}-\{0\}}$
  \task $-x^2\le0\Rightarrow R=\res{(-\infty,\,5]}$
\end{tasks}

\ex{3} El denominador només s'ha d'anul·lar en $x=3$ (arrel doble): $x^2-6x+k=(x-3)^2\Rightarrow\res{k=9}$. El domini és $\mathbb{R}$ si el denominador no s'anul·la mai: $\Delta=36-4k<0\Rightarrow\res{k>9}$.


% =====================================================
\seccio{3}{Punts de tall amb els eixos}

\exf{1} \begin{tasks}(1)
  \task $OY$: $\res{(0,\,4)}$; $OX$: $\res{(1,\,0),\ (4,\,0)}$
  \task $OY$: $\res{(0,\,-6)}$; $OX$: $\res{(3,\,0)}$
  \task $OY$: $\res{(0,\,0)}$; $OX$: $\res{(-3,\,0),\ (0,\,0),\ (3,\,0)}$
  \task $OY$: $\res{(0,\,-2)}$; $OX$: $\res{(-4,\,0)}$
  \task $OY$: $\res{(0,\,1)}$; $OX$: $\res{(-5,\,0)}$
  \task $OY$: $\res{(0,\,5)}$; $OX$: no en té
  \task $OY$: $\res{(0,\,4)}$; $OX$: $\res{(-2,\,0),\ (-1,\,0),\ (1,\,0),\ (2,\,0)}$
  \task $OY$: $\res{(0,\,-4)}$; $OX$: $\res{(-2,\,0),\ (2,\,0)}$
\end{tasks}

\ex{2} $f(2)=4+2a-6=0\Rightarrow\res{a=1}$. \ $x^2+x-6=0\Rightarrow$ l'altre tall és $\res{(-3,\,0)}$. \ $OY$: $\res{(0,\,-6)}$.

\ex{3} Si s'anul·la en $-1$ i en $4$, és $f(x)=a(x+1)(x-4)$. \ $f(0)=-4a=-8\Rightarrow a=2$: \ $\res{f(x)=2(x+1)(x-4)=2x^2-6x-8}$


% =====================================================
\seccio{4}{Creixement i decreixement. Màxims i mínims}

\exf{1} $D=\res{[-5,\,6]}$, $R=\res{[-3,\,4]}$. Creixent en $\res{(-5,\,-3)\cup(-1,\,1)\cup(3,\,6)}$; decreixent en $\res{(-3,\,-1)\cup(1,\,3)}$. Màxims: $\res{(-3,\,3)}$ relatiu i $\res{(1,\,4)}$ absolut. Mínims: $\res{(-1,\,1)}$ relatiu i $\res{(3,\,-3)}$ absolut.

\ex{1} Creixent en $\res{(-\infty,\,-1)\cup(1,\,+\infty)}$; decreixent en $\res{(-1,\,1)}$. Màxim relatiu $\res{(-1,\,2)}$ i mínim relatiu $\res{(1,\,-2)}$. No té extrems absoluts: per la dreta puja sense parar i per l'esquerra baixa sense parar.

\ex{2} Per exemple (n'hi ha moltes), la paràbola $f(x)=\frac12(x-1)^2-2$:

\begin{center}\begin{tikzpicture}[x=0.463cm,y=0.463cm,baseline={(current bounding box.center)},font=\scriptsize]
\draw[muted!25,line width=.3pt] (-3,-3)--(-3,4);
\draw[muted!25,line width=.3pt] (-2,-3)--(-2,4);
\draw[muted!25,line width=.3pt] (-1,-3)--(-1,4);
\draw[muted!25,line width=.3pt] (0,-3)--(0,4);
\draw[muted!25,line width=.3pt] (1,-3)--(1,4);
\draw[muted!25,line width=.3pt] (2,-3)--(2,4);
\draw[muted!25,line width=.3pt] (3,-3)--(3,4);
\draw[muted!25,line width=.3pt] (4,-3)--(4,4);
\draw[muted!25,line width=.3pt] (5,-3)--(5,4);
\draw[muted!25,line width=.3pt] (-3,-3)--(5,-3);
\draw[muted!25,line width=.3pt] (-3,-2)--(5,-2);
\draw[muted!25,line width=.3pt] (-3,-1)--(5,-1);
\draw[muted!25,line width=.3pt] (-3,0)--(5,0);
\draw[muted!25,line width=.3pt] (-3,1)--(5,1);
\draw[muted!25,line width=.3pt] (-3,2)--(5,2);
\draw[muted!25,line width=.3pt] (-3,3)--(5,3);
\draw[muted!25,line width=.3pt] (-3,4)--(5,4);
\draw[-latex,muted] (-3,0)--(5,0);
\draw[-latex,muted] (0,-3)--(0,4);
\node[below=-1pt,text=muted] at (-3,0) {$-3$};
\node[below=-1pt,text=muted] at (-2,0) {$-2$};
\node[below=-1pt,text=muted] at (-1,0) {$-1$};
\node[below=-1pt,text=muted] at (1,0) {$1$};
\node[below=-1pt,text=muted] at (2,0) {$2$};
\node[below=-1pt,text=muted] at (3,0) {$3$};
\node[below=-1pt,text=muted] at (4,0) {$4$};
\node[below=-1pt,text=muted] at (5,0) {$5$};
\node[left=-1pt,text=muted] at (0,-3) {$-3$};
\node[left=-1pt,text=muted] at (0,-2) {$-2$};
\node[left=-1pt,text=muted] at (0,-1) {$-1$};
\node[left=-1pt,text=muted] at (0,1) {$1$};
\node[left=-1pt,text=muted] at (0,2) {$2$};
\node[left=-1pt,text=muted] at (0,3) {$3$};
\node[left=-1pt,text=muted] at (0,4) {$4$};
\begin{scope}\clip (-3,-3) rectangle (5,4);
\draw[accent,line width=1.1pt,line join=round] (-3,6)--(-2.936,5.746)--(-2.872,5.496)--(-2.808,5.25)--(-2.744,5.009)--(-2.68,4.771)--(-2.616,4.538)--(-2.552,4.308)--(-2.488,4.083)--(-2.424,3.862)--(-2.36,3.645)--(-2.296,3.432)--(-2.232,3.223)--(-2.168,3.018)--(-2.104,2.817)--(-2.04,2.621)--(-1.976,2.428)--(-1.912,2.24)--(-1.848,2.056)--(-1.784,1.875)--(-1.72,1.699)--(-1.656,1.527)--(-1.592,1.359)--(-1.528,1.195)--(-1.464,1.036)--(-1.4,0.88)--(-1.336,0.728)--(-1.272,0.581)--(-1.208,0.438)--(-1.144,0.298)--(-1.08,0.163)--(-1.016,0.032)--(-0.952,-0.095)--(-0.888,-0.218)--(-0.824,-0.337)--(-0.76,-0.451)--(-0.696,-0.562)--(-0.632,-0.668)--(-0.568,-0.771)--(-0.504,-0.869)--(-0.44,-0.963)--(-0.376,-1.053)--(-0.312,-1.139)--(-0.248,-1.221)--(-0.184,-1.299)--(-0.12,-1.373)--(-0.056,-1.442)--(0.008,-1.508)--(0.072,-1.569)--(0.136,-1.627)--(0.2,-1.68)--(0.264,-1.729)--(0.328,-1.774)--(0.392,-1.815)--(0.456,-1.852)--(0.52,-1.885)--(0.584,-1.913)--(0.648,-1.938)--(0.712,-1.959)--(0.776,-1.975)--(0.84,-1.987)--(0.904,-1.995)--(0.968,-1.999)--(1.032,-1.999)--(1.096,-1.995)--(1.16,-1.987)--(1.224,-1.975)--(1.288,-1.959)--(1.352,-1.938)--(1.416,-1.913)--(1.48,-1.885)--(1.544,-1.852)--(1.608,-1.815)--(1.672,-1.774)--(1.736,-1.729)--(1.8,-1.68)--(1.864,-1.627)--(1.928,-1.569)--(1.992,-1.508)--(2.056,-1.442)--(2.12,-1.373)--(2.184,-1.299)--(2.248,-1.221)--(2.312,-1.139)--(2.376,-1.053)--(2.44,-0.963)--(2.504,-0.869)--(2.568,-0.771)--(2.632,-0.668)--(2.696,-0.562)--(2.76,-0.451)--(2.824,-0.337)--(2.888,-0.218)--(2.952,-0.095)--(3.016,0.032)--(3.08,0.163)--(3.144,0.298)--(3.208,0.438)--(3.272,0.581)--(3.336,0.728)--(3.4,0.88)--(3.464,1.036)--(3.528,1.195)--(3.592,1.359)--(3.656,1.527)--(3.72,1.699)--(3.784,1.875)--(3.848,2.056)--(3.912,2.24)--(3.976,2.428)--(4.04,2.621)--(4.104,2.817)--(4.168,3.018)--(4.232,3.223)--(4.296,3.432)--(4.36,3.645)--(4.424,3.862)--(4.488,4.083)--(4.552,4.308)--(4.616,4.538)--(4.68,4.771)--(4.744,5.009)--(4.808,5.25)--(4.872,5.496)--(4.936,5.746)--(5,6);
\filldraw[fill=blue!70!black,draw=blue!70!black,line width=.8pt] (1,-2) circle (1.5pt);
\filldraw[fill=blue!70!black,draw=blue!70!black,line width=.8pt] (-1,0) circle (1.5pt);
\filldraw[fill=blue!70!black,draw=blue!70!black,line width=.8pt] (3,0) circle (1.5pt);
\end{scope}\end{tikzpicture}\end{center}


\ex{2} \begin{enumerate}[label=\alph*),leftmargin=*,itemsep=2pt,topsep=2pt]\item $D=\res{[0,\,12]}$: el vol dura 12 min. $R=\res{[0,\,50]}$: l'altura va de 0 m a 50 m.\item Puja en $\res{(0,\,2)\cup(5,\,7)}$, es manté en $\res{(2,\,5)}$ i baixa en $\res{(7,\,12)}$.\item Màxim absolut $\res{(7,\,50)}$: 50 m al cap de 7 min.\item Des del minut 5 fins al minut 9 (de 7 a 10 baixa 10 m cada minut): \rest{4 minuts}.\end{enumerate}


% =====================================================
\seccio{5}{Continuïtat}

\exf{1} $x=-3$: $\res{\text{evitable}}$ (el punt està fora de lloc). \ $x=0$: $\res{\text{salt finit}}$ (de 3 a 1). \ $x=2$: $\res{\text{salt infinit}}$ (per l'esquerra arriba a 0 i per la dreta se'n va a $+\infty$).

\ex{2} $x=-3$: $\res{\text{evitable}}$ (els dos costats arriben a 2, però $f(-3)=4$). \ $x=1$: $\res{\text{salt infinit}}$ (per l'esquerra arriba a 0 i per la dreta se'n va a $+\infty$). \ $x=3$: $\res{\text{salt finit}}$ (per l'esquerra arriba a 1 i $f(3)=2$). És contínua en $\res{[-6,\,-3)\cup(-3,\,1)\cup(1,\,3)\cup[3,\,7]}$.

\ex{2} \begin{tasks}(1)
  \task Discontinuïtat $\res{\text{evitable}}$: els dos costats arriben a 3, però $f(2)$ no existeix. Seria contínua si $\res{f(2)=3}$.
  \task Discontinuïtat de $\res{\text{salt finit}}$: per l'esquerra arriba a 1 (i $f(2)=1$) i per la dreta a 3.
  \task $\res{\text{Contínua}}$: els dos costats arriben a 1 i $f(2)=1$. Té un pic, però es dibuixa sense alçar el llapis.
  \task Discontinuïtat $\res{\text{asimptòtica}}$: per l'esquerra se'n va a $-\infty$ i per la dreta a $+\infty$ (asímptota vertical $x=2$).
\end{tasks}

\ex{2} Per exemple (n'hi ha moltes):

\begin{center}\begin{tikzpicture}[x=0.46cm,y=0.46cm,baseline={(current bounding box.center)},font=\scriptsize]
\draw[muted!25,line width=.3pt] (-3,-2)--(-3,6);
\draw[muted!25,line width=.3pt] (-2,-2)--(-2,6);
\draw[muted!25,line width=.3pt] (-1,-2)--(-1,6);
\draw[muted!25,line width=.3pt] (0,-2)--(0,6);
\draw[muted!25,line width=.3pt] (1,-2)--(1,6);
\draw[muted!25,line width=.3pt] (2,-2)--(2,6);
\draw[muted!25,line width=.3pt] (3,-2)--(3,6);
\draw[muted!25,line width=.3pt] (4,-2)--(4,6);
\draw[muted!25,line width=.3pt] (5,-2)--(5,6);
\draw[muted!25,line width=.3pt] (6,-2)--(6,6);
\draw[muted!25,line width=.3pt] (7,-2)--(7,6);
\draw[muted!25,line width=.3pt] (-3,-2)--(7,-2);
\draw[muted!25,line width=.3pt] (-3,-1)--(7,-1);
\draw[muted!25,line width=.3pt] (-3,0)--(7,0);
\draw[muted!25,line width=.3pt] (-3,1)--(7,1);
\draw[muted!25,line width=.3pt] (-3,2)--(7,2);
\draw[muted!25,line width=.3pt] (-3,3)--(7,3);
\draw[muted!25,line width=.3pt] (-3,4)--(7,4);
\draw[muted!25,line width=.3pt] (-3,5)--(7,5);
\draw[muted!25,line width=.3pt] (-3,6)--(7,6);
\draw[-latex,muted] (-3,0)--(7,0);
\draw[-latex,muted] (0,-2)--(0,6);
\node[below=-1pt,text=muted] at (-3,0) {$-3$};
\node[below=-1pt,text=muted] at (-2,0) {$-2$};
\node[below=-1pt,text=muted] at (-1,0) {$-1$};
\node[below=-1pt,text=muted] at (1,0) {$1$};
\node[below=-1pt,text=muted] at (2,0) {$2$};
\node[below=-1pt,text=muted] at (3,0) {$3$};
\node[below=-1pt,text=muted] at (4,0) {$4$};
\node[below=-1pt,text=muted] at (5,0) {$5$};
\node[below=-1pt,text=muted] at (6,0) {$6$};
\node[below=-1pt,text=muted] at (7,0) {$7$};
\node[left=-1pt,text=muted] at (0,-2) {$-2$};
\node[left=-1pt,text=muted] at (0,-1) {$-1$};
\node[left=-1pt,text=muted] at (0,1) {$1$};
\node[left=-1pt,text=muted] at (0,2) {$2$};
\node[left=-1pt,text=muted] at (0,3) {$3$};
\node[left=-1pt,text=muted] at (0,4) {$4$};
\node[left=-1pt,text=muted] at (0,5) {$5$};
\node[left=-1pt,text=muted] at (0,6) {$6$};
\begin{scope}\clip (-3,-2) rectangle (7,6);
\draw[accent,line width=1.1pt,line join=round] (-3,-2)--(-2.968,-1.968)--(-2.936,-1.936)--(-2.904,-1.904)--(-2.872,-1.872)--(-2.84,-1.84)--(-2.808,-1.808)--(-2.776,-1.776)--(-2.744,-1.744)--(-2.712,-1.712)--(-2.68,-1.68)--(-2.648,-1.648)--(-2.616,-1.616)--(-2.584,-1.584)--(-2.552,-1.552)--(-2.52,-1.52)--(-2.488,-1.488)--(-2.456,-1.456)--(-2.424,-1.424)--(-2.392,-1.392)--(-2.36,-1.36)--(-2.328,-1.328)--(-2.296,-1.296)--(-2.264,-1.264)--(-2.232,-1.232)--(-2.2,-1.2)--(-2.168,-1.168)--(-2.136,-1.136)--(-2.104,-1.104)--(-2.072,-1.072)--(-2.04,-1.04)--(-2.008,-1.008)--(-1.976,-0.976)--(-1.944,-0.944)--(-1.912,-0.912)--(-1.88,-0.88)--(-1.848,-0.848)--(-1.816,-0.816)--(-1.784,-0.784)--(-1.752,-0.752)--(-1.72,-0.72)--(-1.688,-0.688)--(-1.656,-0.656)--(-1.624,-0.624)--(-1.592,-0.592)--(-1.56,-0.56)--(-1.528,-0.528)--(-1.496,-0.496)--(-1.464,-0.464)--(-1.432,-0.432)--(-1.4,-0.4)--(-1.368,-0.368)--(-1.336,-0.336)--(-1.304,-0.304)--(-1.272,-0.272)--(-1.24,-0.24)--(-1.208,-0.208)--(-1.176,-0.176)--(-1.144,-0.144)--(-1.112,-0.112)--(-1.08,-0.08)--(-1.048,-0.048)--(-1.016,-0.016)--(-0.984,0.016)--(-0.952,0.048)--(-0.92,0.08)--(-0.888,0.112)--(-0.856,0.144)--(-0.824,0.176)--(-0.792,0.208)--(-0.76,0.24)--(-0.728,0.272)--(-0.696,0.304)--(-0.664,0.336)--(-0.632,0.368)--(-0.6,0.4)--(-0.568,0.432)--(-0.536,0.464)--(-0.504,0.496)--(-0.472,0.528)--(-0.44,0.56)--(-0.408,0.592)--(-0.376,0.624)--(-0.344,0.656)--(-0.312,0.688)--(-0.28,0.72)--(-0.248,0.752)--(-0.216,0.784)--(-0.184,0.816)--(-0.152,0.848)--(-0.12,0.88)--(-0.088,0.912)--(-0.056,0.944)--(-0.024,0.976)--(0.008,1.008)--(0.04,1.04)--(0.072,1.072)--(0.104,1.104)--(0.136,1.136)--(0.168,1.168)--(0.2,1.2)--(0.232,1.232)--(0.264,1.264)--(0.296,1.296)--(0.328,1.328)--(0.36,1.36)--(0.392,1.392)--(0.424,1.424)--(0.456,1.456)--(0.488,1.488)--(0.52,1.52)--(0.552,1.552)--(0.584,1.584)--(0.616,1.616)--(0.648,1.648)--(0.68,1.68)--(0.712,1.712)--(0.744,1.744)--(0.776,1.776)--(0.808,1.808)--(0.84,1.84)--(0.872,1.872)--(0.904,1.904)--(0.936,1.936)--(0.968,1.968)--(1,2);
\draw[accent,line width=1.1pt,line join=round] (1,0)--(1.048,0.048)--(1.096,0.096)--(1.144,0.144)--(1.192,0.192)--(1.24,0.24)--(1.288,0.288)--(1.336,0.336)--(1.384,0.384)--(1.432,0.432)--(1.48,0.48)--(1.528,0.528)--(1.576,0.576)--(1.624,0.624)--(1.672,0.672)--(1.72,0.72)--(1.768,0.768)--(1.816,0.816)--(1.864,0.864)--(1.912,0.912)--(1.96,0.96)--(2.008,1.008)--(2.056,1.056)--(2.104,1.104)--(2.152,1.152)--(2.2,1.2)--(2.248,1.248)--(2.296,1.296)--(2.344,1.344)--(2.392,1.392)--(2.44,1.44)--(2.488,1.488)--(2.536,1.536)--(2.584,1.584)--(2.632,1.632)--(2.68,1.68)--(2.728,1.728)--(2.776,1.776)--(2.824,1.824)--(2.872,1.872)--(2.92,1.92)--(2.968,1.968)--(3.016,2.016)--(3.064,2.064)--(3.112,2.112)--(3.16,2.16)--(3.208,2.208)--(3.256,2.256)--(3.304,2.304)--(3.352,2.352)--(3.4,2.4)--(3.448,2.448)--(3.496,2.496)--(3.544,2.544)--(3.592,2.592)--(3.64,2.64)--(3.688,2.688)--(3.736,2.736)--(3.784,2.784)--(3.832,2.832)--(3.88,2.88)--(3.928,2.928)--(3.976,2.976)--(4.024,3.024)--(4.072,3.072)--(4.12,3.12)--(4.168,3.168)--(4.216,3.216)--(4.264,3.264)--(4.312,3.312)--(4.36,3.36)--(4.408,3.408)--(4.456,3.456)--(4.504,3.504)--(4.552,3.552)--(4.6,3.6)--(4.648,3.648)--(4.696,3.696)--(4.744,3.744)--(4.792,3.792)--(4.84,3.84)--(4.888,3.888)--(4.936,3.936)--(4.984,3.984)--(5.032,4.032)--(5.08,4.08)--(5.128,4.128)--(5.176,4.176)--(5.224,4.224)--(5.272,4.272)--(5.32,4.32)--(5.368,4.368)--(5.416,4.416)--(5.464,4.464)--(5.512,4.512)--(5.56,4.56)--(5.608,4.608)--(5.656,4.656)--(5.704,4.704)--(5.752,4.752)--(5.8,4.8)--(5.848,4.848)--(5.896,4.896)--(5.944,4.944)--(5.992,4.992)--(6.04,5.04)--(6.088,5.088)--(6.136,5.136)--(6.184,5.184)--(6.232,5.232)--(6.28,5.28)--(6.328,5.328)--(6.376,5.376)--(6.424,5.424)--(6.472,5.472)--(6.52,5.52)--(6.568,5.568)--(6.616,5.616)--(6.664,5.664)--(6.712,5.712)--(6.76,5.76)--(6.808,5.808)--(6.856,5.856)--(6.904,5.904)--(6.952,5.952)--(7,6);
\filldraw[fill=accent,draw=accent,line width=.8pt] (1,2) circle (1.5pt);
\filldraw[fill=white,draw=accent,line width=.8pt] (1,0) circle (1.5pt);
\filldraw[fill=white,draw=accent,line width=.8pt] (4,3) circle (1.5pt);
\filldraw[fill=accent,draw=accent,line width=.8pt] (4,5) circle (1.5pt);
\end{scope}\end{tikzpicture}\end{center}



% =====================================================
\seccio{6}{Simetries: funcions parelles i senars}

\exf{1} \begin{tasks}(4)
  \task \rest{Parella}
  \task \rest{Senar}
  \task \rest{Cap}
  \task \rest{Senar}
  \task \rest{Parella}
  \task \rest{Parella}
  \task \rest{Cap}
  \task \rest{Parella}
\end{tasks}

\ex{1} \begin{tasks}(4)
  \task \rest{Parella}
  \task \rest{Senar}
  \task \rest{Cap}
  \task \rest{Senar}
\end{tasks}

\ex{2} \begin{tasks}(2)
  \task Parella (simetria respecte de $OY$):

\begin{center}\begin{tikzpicture}[x=0.36cm,y=0.372cm,baseline={(current bounding box.center)},font=\scriptsize]
\draw[muted!25,line width=.3pt] (-5,-3)--(-5,3);
\draw[muted!25,line width=.3pt] (-4,-3)--(-4,3);
\draw[muted!25,line width=.3pt] (-3,-3)--(-3,3);
\draw[muted!25,line width=.3pt] (-2,-3)--(-2,3);
\draw[muted!25,line width=.3pt] (-1,-3)--(-1,3);
\draw[muted!25,line width=.3pt] (0,-3)--(0,3);
\draw[muted!25,line width=.3pt] (1,-3)--(1,3);
\draw[muted!25,line width=.3pt] (2,-3)--(2,3);
\draw[muted!25,line width=.3pt] (3,-3)--(3,3);
\draw[muted!25,line width=.3pt] (4,-3)--(4,3);
\draw[muted!25,line width=.3pt] (5,-3)--(5,3);
\draw[muted!25,line width=.3pt] (-5,-3)--(5,-3);
\draw[muted!25,line width=.3pt] (-5,-2)--(5,-2);
\draw[muted!25,line width=.3pt] (-5,-1)--(5,-1);
\draw[muted!25,line width=.3pt] (-5,0)--(5,0);
\draw[muted!25,line width=.3pt] (-5,1)--(5,1);
\draw[muted!25,line width=.3pt] (-5,2)--(5,2);
\draw[muted!25,line width=.3pt] (-5,3)--(5,3);
\draw[-latex,muted] (-5,0)--(5,0);
\draw[-latex,muted] (0,-3)--(0,3);
\node[below=-1pt,text=muted] at (-5,0) {$-5$};
\node[below=-1pt,text=muted] at (-4,0) {$-4$};
\node[below=-1pt,text=muted] at (-3,0) {$-3$};
\node[below=-1pt,text=muted] at (-2,0) {$-2$};
\node[below=-1pt,text=muted] at (-1,0) {$-1$};
\node[below=-1pt,text=muted] at (1,0) {$1$};
\node[below=-1pt,text=muted] at (2,0) {$2$};
\node[below=-1pt,text=muted] at (3,0) {$3$};
\node[below=-1pt,text=muted] at (4,0) {$4$};
\node[below=-1pt,text=muted] at (5,0) {$5$};
\node[left=-1pt,text=muted] at (0,-3) {$-3$};
\node[left=-1pt,text=muted] at (0,-2) {$-2$};
\node[left=-1pt,text=muted] at (0,-1) {$-1$};
\node[left=-1pt,text=muted] at (0,1) {$1$};
\node[left=-1pt,text=muted] at (0,2) {$2$};
\node[left=-1pt,text=muted] at (0,3) {$3$};
\begin{scope}\clip (-5,-3) rectangle (5,3);
\draw[accent,line width=1.1pt,line join=round] (0,0)--(1,2);
\draw[accent,line width=1.1pt,line join=round] (1,2)--(3,-1);
\draw[accent,line width=1.1pt,line join=round] (3,-1)--(4,0);
\draw[blue!70!black,line width=1.1pt,line join=round] (0,0)--(-1,2);
\draw[blue!70!black,line width=1.1pt,line join=round] (-1,2)--(-3,-1);
\draw[blue!70!black,line width=1.1pt,line join=round] (-3,-1)--(-4,0);
\end{scope}\end{tikzpicture}\end{center}

  \task Senar (simetria respecte de l'origen):

\begin{center}\begin{tikzpicture}[x=0.36cm,y=0.372cm,baseline={(current bounding box.center)},font=\scriptsize]
\draw[muted!25,line width=.3pt] (-5,-3)--(-5,3);
\draw[muted!25,line width=.3pt] (-4,-3)--(-4,3);
\draw[muted!25,line width=.3pt] (-3,-3)--(-3,3);
\draw[muted!25,line width=.3pt] (-2,-3)--(-2,3);
\draw[muted!25,line width=.3pt] (-1,-3)--(-1,3);
\draw[muted!25,line width=.3pt] (0,-3)--(0,3);
\draw[muted!25,line width=.3pt] (1,-3)--(1,3);
\draw[muted!25,line width=.3pt] (2,-3)--(2,3);
\draw[muted!25,line width=.3pt] (3,-3)--(3,3);
\draw[muted!25,line width=.3pt] (4,-3)--(4,3);
\draw[muted!25,line width=.3pt] (5,-3)--(5,3);
\draw[muted!25,line width=.3pt] (-5,-3)--(5,-3);
\draw[muted!25,line width=.3pt] (-5,-2)--(5,-2);
\draw[muted!25,line width=.3pt] (-5,-1)--(5,-1);
\draw[muted!25,line width=.3pt] (-5,0)--(5,0);
\draw[muted!25,line width=.3pt] (-5,1)--(5,1);
\draw[muted!25,line width=.3pt] (-5,2)--(5,2);
\draw[muted!25,line width=.3pt] (-5,3)--(5,3);
\draw[-latex,muted] (-5,0)--(5,0);
\draw[-latex,muted] (0,-3)--(0,3);
\node[below=-1pt,text=muted] at (-5,0) {$-5$};
\node[below=-1pt,text=muted] at (-4,0) {$-4$};
\node[below=-1pt,text=muted] at (-3,0) {$-3$};
\node[below=-1pt,text=muted] at (-2,0) {$-2$};
\node[below=-1pt,text=muted] at (-1,0) {$-1$};
\node[below=-1pt,text=muted] at (1,0) {$1$};
\node[below=-1pt,text=muted] at (2,0) {$2$};
\node[below=-1pt,text=muted] at (3,0) {$3$};
\node[below=-1pt,text=muted] at (4,0) {$4$};
\node[below=-1pt,text=muted] at (5,0) {$5$};
\node[left=-1pt,text=muted] at (0,-3) {$-3$};
\node[left=-1pt,text=muted] at (0,-2) {$-2$};
\node[left=-1pt,text=muted] at (0,-1) {$-1$};
\node[left=-1pt,text=muted] at (0,1) {$1$};
\node[left=-1pt,text=muted] at (0,2) {$2$};
\node[left=-1pt,text=muted] at (0,3) {$3$};
\begin{scope}\clip (-5,-3) rectangle (5,3);
\draw[accent,line width=1.1pt,line join=round] (0,0)--(1,2);
\draw[accent,line width=1.1pt,line join=round] (1,2)--(3,-1);
\draw[accent,line width=1.1pt,line join=round] (3,-1)--(4,0);
\draw[blue!70!black,line width=1.1pt,line join=round] (0,0)--(-1,-2);
\draw[blue!70!black,line width=1.1pt,line join=round] (-1,-2)--(-3,1);
\draw[blue!70!black,line width=1.1pt,line join=round] (-3,1)--(-4,0);
\end{scope}\end{tikzpicture}\end{center}

\end{tasks}

\ex{3} $(f+g)(-x)=f(-x)+g(-x)=-f(x)-g(x)=-(f+g)(x)$: \rest{senar}. \ $(f\cdot g)(-x)=\big(-f(x)\big)\cdot\big(-g(x)\big)=f(x)\cdot g(x)$: \rest{parella}. \ Amb l'exemple: $x^3+x$ és senar i $x^3\cdot x=x^4$ és parella.


% =====================================================
\seccio{7}{Periodicitat}

\exf{1} \begin{tasks}(1)
  \task $\res{T=3}$
  \task $f(10)=f(1+3\cdot3)=f(1)=\res{2}$; \ $f(-1)=f(-1+3)=f(2)=\res{3}$; \ $f(30)=f(0+10\cdot3)=f(0)=\res{1}$; \ $f(47)=f(2+15\cdot3)=f(2)=\res{3}$
\end{tasks}

\ex{2} $f(11)=f(1+2\cdot5)=\res{4}$; \ $f(-2)=f(-2+5)=f(3)=\res{-2}$; \ $f(23)=f(3+4\cdot5)=\res{-2}$; \ $f(101)=f(1+20\cdot5)=\res{4}$

\ex{2} \begin{enumerate}[label=\alph*),leftmargin=*,itemsep=2pt,topsep=2pt]\item $\res{T=60\text{ s}}$. Altura mínima 2 m i màxima $2+20=22$ m: $R=\res{[2,\,22]}$.\item $150=2\cdot60+30$: com als 30 s (mitja volta), és dalt de tot: \rest{22 m}. \ $240=4\cdot60$: com a l'inici: \rest{2 m}.\item Fa $600:60=10$ voltes: \rest{10 vegades}.\end{enumerate}

\ex{2} 

\begin{center}\begin{tikzpicture}[x=0.571cm,y=1.24cm,baseline={(current bounding box.center)},font=\scriptsize]
\draw[muted!25,line width=.3pt] (-5,-1)--(-5,3);
\draw[muted!25,line width=.3pt] (-4,-1)--(-4,3);
\draw[muted!25,line width=.3pt] (-3,-1)--(-3,3);
\draw[muted!25,line width=.3pt] (-2,-1)--(-2,3);
\draw[muted!25,line width=.3pt] (-1,-1)--(-1,3);
\draw[muted!25,line width=.3pt] (0,-1)--(0,3);
\draw[muted!25,line width=.3pt] (1,-1)--(1,3);
\draw[muted!25,line width=.3pt] (2,-1)--(2,3);
\draw[muted!25,line width=.3pt] (3,-1)--(3,3);
\draw[muted!25,line width=.3pt] (4,-1)--(4,3);
\draw[muted!25,line width=.3pt] (5,-1)--(5,3);
\draw[muted!25,line width=.3pt] (6,-1)--(6,3);
\draw[muted!25,line width=.3pt] (7,-1)--(7,3);
\draw[muted!25,line width=.3pt] (8,-1)--(8,3);
\draw[muted!25,line width=.3pt] (9,-1)--(9,3);
\draw[muted!25,line width=.3pt] (-5,-1)--(9,-1);
\draw[muted!25,line width=.3pt] (-5,0)--(9,0);
\draw[muted!25,line width=.3pt] (-5,1)--(9,1);
\draw[muted!25,line width=.3pt] (-5,2)--(9,2);
\draw[muted!25,line width=.3pt] (-5,3)--(9,3);
\draw[-latex,muted] (-5,0)--(9,0);
\draw[-latex,muted] (0,-1)--(0,3);
\node[below=-1pt,text=muted] at (-5,0) {$-5$};
\node[below=-1pt,text=muted] at (-4,0) {$-4$};
\node[below=-1pt,text=muted] at (-3,0) {$-3$};
\node[below=-1pt,text=muted] at (-2,0) {$-2$};
\node[below=-1pt,text=muted] at (-1,0) {$-1$};
\node[below=-1pt,text=muted] at (1,0) {$1$};
\node[below=-1pt,text=muted] at (2,0) {$2$};
\node[below=-1pt,text=muted] at (3,0) {$3$};
\node[below=-1pt,text=muted] at (4,0) {$4$};
\node[below=-1pt,text=muted] at (5,0) {$5$};
\node[below=-1pt,text=muted] at (6,0) {$6$};
\node[below=-1pt,text=muted] at (7,0) {$7$};
\node[below=-1pt,text=muted] at (8,0) {$8$};
\node[below=-1pt,text=muted] at (9,0) {$9$};
\node[left=-1pt,text=muted] at (0,-1) {$-1$};
\node[left=-1pt,text=muted] at (0,1) {$1$};
\node[left=-1pt,text=muted] at (0,2) {$2$};
\node[left=-1pt,text=muted] at (0,3) {$3$};
\begin{scope}\clip (-5,-1) rectangle (9,3);
\draw[accent,line width=1.1pt,line join=round] (-4,2)--(-3.9,1.9)--(-3.8,1.8)--(-3.7,1.7)--(-3.6,1.6)--(-3.5,1.5)--(-3.4,1.4)--(-3.3,1.3)--(-3.2,1.2)--(-3.1,1.1)--(-3,1)--(-2.9,0.9)--(-2.8,0.8)--(-2.7,0.7)--(-2.6,0.6)--(-2.5,0.5)--(-2.4,0.4)--(-2.3,0.3)--(-2.2,0.2)--(-2.1,0.1)--(-2,0)--(-1.9,0.1)--(-1.8,0.2)--(-1.7,0.3)--(-1.6,0.4)--(-1.5,0.5)--(-1.4,0.6)--(-1.3,0.7)--(-1.2,0.8)--(-1.1,0.9)--(-1,1)--(-0.9,1.1)--(-0.8,1.2)--(-0.7,1.3)--(-0.6,1.4)--(-0.5,1.5)--(-0.4,1.6)--(-0.3,1.7)--(-0.2,1.8)--(-0.1,1.9)--(0,2)--(0.1,1.9)--(0.2,1.8)--(0.3,1.7)--(0.4,1.6)--(0.5,1.5)--(0.6,1.4)--(0.7,1.3)--(0.8,1.2)--(0.9,1.1)--(1,1)--(1.1,0.9)--(1.2,0.8)--(1.3,0.7)--(1.4,0.6)--(1.5,0.5)--(1.6,0.4)--(1.7,0.3)--(1.8,0.2)--(1.9,0.1)--(2,0)--(2.1,0.1)--(2.2,0.2)--(2.3,0.3)--(2.4,0.4)--(2.5,0.5)--(2.6,0.6)--(2.7,0.7)--(2.8,0.8)--(2.9,0.9)--(3,1)--(3.1,1.1)--(3.2,1.2)--(3.3,1.3)--(3.4,1.4)--(3.5,1.5)--(3.6,1.6)--(3.7,1.7)--(3.8,1.8)--(3.9,1.9)--(4,2)--(4.1,1.9)--(4.2,1.8)--(4.3,1.7)--(4.4,1.6)--(4.5,1.5)--(4.6,1.4)--(4.7,1.3)--(4.8,1.2)--(4.9,1.1)--(5,1)--(5.1,0.9)--(5.2,0.8)--(5.3,0.7)--(5.4,0.6)--(5.5,0.5)--(5.6,0.4)--(5.7,0.3)--(5.8,0.2)--(5.9,0.1)--(6,0)--(6.1,0.1)--(6.2,0.2)--(6.3,0.3)--(6.4,0.4)--(6.5,0.5)--(6.6,0.6)--(6.7,0.7)--(6.8,0.8)--(6.9,0.9)--(7,1)--(7.1,1.1)--(7.2,1.2)--(7.3,1.3)--(7.4,1.4)--(7.5,1.5)--(7.6,1.6)--(7.7,1.7)--(7.8,1.8)--(7.9,1.9)--(8,2);
\draw[blue!70!black,line width=1.1pt,line join=round] (0,2)--(0.03,1.97)--(0.06,1.94)--(0.09,1.91)--(0.12,1.88)--(0.15,1.85)--(0.18,1.82)--(0.21,1.79)--(0.24,1.76)--(0.27,1.73)--(0.3,1.7)--(0.33,1.67)--(0.36,1.64)--(0.39,1.61)--(0.42,1.58)--(0.45,1.55)--(0.48,1.52)--(0.51,1.49)--(0.54,1.46)--(0.57,1.43)--(0.6,1.4)--(0.63,1.37)--(0.66,1.34)--(0.69,1.31)--(0.72,1.28)--(0.75,1.25)--(0.78,1.22)--(0.81,1.19)--(0.84,1.16)--(0.87,1.13)--(0.9,1.1)--(0.93,1.07)--(0.96,1.04)--(0.99,1.01)--(1.02,0.98)--(1.05,0.95)--(1.08,0.92)--(1.11,0.89)--(1.14,0.86)--(1.17,0.83)--(1.2,0.8)--(1.23,0.77)--(1.26,0.74)--(1.29,0.71)--(1.32,0.68)--(1.35,0.65)--(1.38,0.62)--(1.41,0.59)--(1.44,0.56)--(1.47,0.53)--(1.5,0.5)--(1.53,0.47)--(1.56,0.44)--(1.59,0.41)--(1.62,0.38)--(1.65,0.35)--(1.68,0.32)--(1.71,0.29)--(1.74,0.26)--(1.77,0.23)--(1.8,0.2)--(1.83,0.17)--(1.86,0.14)--(1.89,0.11)--(1.92,0.08)--(1.95,0.05)--(1.98,0.02)--(2.01,0.01)--(2.04,0.04)--(2.07,0.07)--(2.1,0.1)--(2.13,0.13)--(2.16,0.16)--(2.19,0.19)--(2.22,0.22)--(2.25,0.25)--(2.28,0.28)--(2.31,0.31)--(2.34,0.34)--(2.37,0.37)--(2.4,0.4)--(2.43,0.43)--(2.46,0.46)--(2.49,0.49)--(2.52,0.52)--(2.55,0.55)--(2.58,0.58)--(2.61,0.61)--(2.64,0.64)--(2.67,0.67)--(2.7,0.7)--(2.73,0.73)--(2.76,0.76)--(2.79,0.79)--(2.82,0.82)--(2.85,0.85)--(2.88,0.88)--(2.91,0.91)--(2.94,0.94)--(2.97,0.97)--(3,1)--(3.03,1.03)--(3.06,1.06)--(3.09,1.09)--(3.12,1.12)--(3.15,1.15)--(3.18,1.18)--(3.21,1.21)--(3.24,1.24)--(3.27,1.27)--(3.3,1.3)--(3.33,1.33)--(3.36,1.36)--(3.39,1.39)--(3.42,1.42)--(3.45,1.45)--(3.48,1.48)--(3.51,1.51)--(3.54,1.54)--(3.57,1.57)--(3.6,1.6)--(3.63,1.63)--(3.66,1.66)--(3.69,1.69)--(3.72,1.72)--(3.75,1.75)--(3.78,1.78)--(3.81,1.81)--(3.84,1.84)--(3.87,1.87)--(3.9,1.9)--(3.93,1.93)--(3.96,1.96)--(3.99,1.99)--(4,2);
\end{scope}\end{tikzpicture}\end{center}
$f(13)=f(1+3\cdot4)=f(1)=|1-2|=\res{1}$


% =====================================================
\seccio{8}{Funció inversa}

\exf{1} $f^{-1}(5)=\res{2}$, \ $f^{-1}(-1)=\res{4}$, \ $f^{-1}\big(f(2)\big)=\res{2}$. La gràfica de $f^{-1}$ passa per $\res{(3,\,1)}$, $\res{(5,\,2)}$, $\res{(-1,\,4)}$: s'intercanvien les coordenades.

\ex{2} \begin{tasks}[after-item-skip=8pt](3)
  \task $f^{-1}(x)=\res{\dfrac{x+2}{5}}$
  \task $f^{-1}(x)=\res{\dfrac{2x-1}{3}}$
  \task $f^{-1}(x)=\res{\sqrt[3]{x+4}}$
  \task $f^{-1}(x)=\res{\dfrac4x}$
  \task $f^{-1}(x)=\res{\dfrac{x+2}{x-1}}$
  \task $f^{-1}(x)=\res{x^2+1\ \ (x\geq0)}$
\end{tasks}Les funcions d) i e) són la seua pròpia inversa. En e): $x(y-1)=y+2\Rightarrow xy-y=x+2\Rightarrow y=\frac{x+2}{x-1}$.

\ex{2} $f\big(g(x)\big)=2\left(\frac x2+3\right)-6=x+6-6=\res{x}$; \ $g\big(f(x)\big)=\frac{2x-6}{2}+3=x-3+3=\res{x}$. Les gràfiques són simètriques respecte de la recta $y=x$ i es tallen sobre ella, en $(6,\,6)$.

\begin{center}\begin{tikzpicture}[x=0.323cm,y=0.323cm,baseline={(current bounding box.center)},font=\scriptsize]
\draw[muted!25,line width=.3pt] (-4,-4)--(-4,9);
\draw[muted!25,line width=.3pt] (-3,-4)--(-3,9);
\draw[muted!25,line width=.3pt] (-2,-4)--(-2,9);
\draw[muted!25,line width=.3pt] (-1,-4)--(-1,9);
\draw[muted!25,line width=.3pt] (0,-4)--(0,9);
\draw[muted!25,line width=.3pt] (1,-4)--(1,9);
\draw[muted!25,line width=.3pt] (2,-4)--(2,9);
\draw[muted!25,line width=.3pt] (3,-4)--(3,9);
\draw[muted!25,line width=.3pt] (4,-4)--(4,9);
\draw[muted!25,line width=.3pt] (5,-4)--(5,9);
\draw[muted!25,line width=.3pt] (6,-4)--(6,9);
\draw[muted!25,line width=.3pt] (7,-4)--(7,9);
\draw[muted!25,line width=.3pt] (8,-4)--(8,9);
\draw[muted!25,line width=.3pt] (9,-4)--(9,9);
\draw[muted!25,line width=.3pt] (-4,-4)--(9,-4);
\draw[muted!25,line width=.3pt] (-4,-3)--(9,-3);
\draw[muted!25,line width=.3pt] (-4,-2)--(9,-2);
\draw[muted!25,line width=.3pt] (-4,-1)--(9,-1);
\draw[muted!25,line width=.3pt] (-4,0)--(9,0);
\draw[muted!25,line width=.3pt] (-4,1)--(9,1);
\draw[muted!25,line width=.3pt] (-4,2)--(9,2);
\draw[muted!25,line width=.3pt] (-4,3)--(9,3);
\draw[muted!25,line width=.3pt] (-4,4)--(9,4);
\draw[muted!25,line width=.3pt] (-4,5)--(9,5);
\draw[muted!25,line width=.3pt] (-4,6)--(9,6);
\draw[muted!25,line width=.3pt] (-4,7)--(9,7);
\draw[muted!25,line width=.3pt] (-4,8)--(9,8);
\draw[muted!25,line width=.3pt] (-4,9)--(9,9);
\draw[-latex,muted] (-4,0)--(9,0);
\draw[-latex,muted] (0,-4)--(0,9);
\node[below=-1pt,text=muted] at (-4,0) {$-4$};
\node[below=-1pt,text=muted] at (-2,0) {$-2$};
\node[below=-1pt,text=muted] at (2,0) {$2$};
\node[below=-1pt,text=muted] at (4,0) {$4$};
\node[below=-1pt,text=muted] at (6,0) {$6$};
\node[below=-1pt,text=muted] at (8,0) {$8$};
\node[left=-1pt,text=muted] at (0,-4) {$-4$};
\node[left=-1pt,text=muted] at (0,-2) {$-2$};
\node[left=-1pt,text=muted] at (0,2) {$2$};
\node[left=-1pt,text=muted] at (0,4) {$4$};
\node[left=-1pt,text=muted] at (0,6) {$6$};
\node[left=-1pt,text=muted] at (0,8) {$8$};
\begin{scope}\clip (-4,-4) rectangle (9,9);
\draw[muted,line width=1.1pt,dashed,line join=round] (-4,-4)--(9,9);
\draw[accent,line width=1.1pt,line join=round] (-4,-14)--(-3.896,-13.792)--(-3.792,-13.584)--(-3.688,-13.376)--(-3.584,-13.168)--(-3.48,-12.96)--(-3.376,-12.752)--(-3.272,-12.544)--(-3.168,-12.336)--(-3.064,-12.128)--(-2.96,-11.92)--(-2.856,-11.712)--(-2.752,-11.504)--(-2.648,-11.296)--(-2.544,-11.088)--(-2.44,-10.88)--(-2.336,-10.672)--(-2.232,-10.464)--(-2.128,-10.256)--(-2.024,-10.048)--(-1.92,-9.84)--(-1.816,-9.632)--(-1.712,-9.424)--(-1.608,-9.216)--(-1.504,-9.008)--(-1.4,-8.8)--(-1.296,-8.592)--(-1.192,-8.384)--(-1.088,-8.176)--(-0.984,-7.968)--(-0.88,-7.76)--(-0.776,-7.552)--(-0.672,-7.344)--(-0.568,-7.136)--(-0.464,-6.928)--(-0.36,-6.72)--(-0.256,-6.512)--(-0.152,-6.304)--(-0.048,-6.096)--(0.056,-5.888)--(0.16,-5.68)--(0.264,-5.472)--(0.368,-5.264)--(0.472,-5.056)--(0.576,-4.848)--(0.68,-4.64)--(0.784,-4.432)--(0.888,-4.224)--(0.992,-4.016)--(1.096,-3.808)--(1.2,-3.6)--(1.304,-3.392)--(1.408,-3.184)--(1.512,-2.976)--(1.616,-2.768)--(1.72,-2.56)--(1.824,-2.352)--(1.928,-2.144)--(2.032,-1.936)--(2.136,-1.728)--(2.24,-1.52)--(2.344,-1.312)--(2.448,-1.104)--(2.552,-0.896)--(2.656,-0.688)--(2.76,-0.48)--(2.864,-0.272)--(2.968,-0.064)--(3.072,0.144)--(3.176,0.352)--(3.28,0.56)--(3.384,0.768)--(3.488,0.976)--(3.592,1.184)--(3.696,1.392)--(3.8,1.6)--(3.904,1.808)--(4.008,2.016)--(4.112,2.224)--(4.216,2.432)--(4.32,2.64)--(4.424,2.848)--(4.528,3.056)--(4.632,3.264)--(4.736,3.472)--(4.84,3.68)--(4.944,3.888)--(5.048,4.096)--(5.152,4.304)--(5.256,4.512)--(5.36,4.72)--(5.464,4.928)--(5.568,5.136)--(5.672,5.344)--(5.776,5.552)--(5.88,5.76)--(5.984,5.968)--(6.088,6.176)--(6.192,6.384)--(6.296,6.592)--(6.4,6.8)--(6.504,7.008)--(6.608,7.216)--(6.712,7.424)--(6.816,7.632)--(6.92,7.84)--(7.024,8.048)--(7.128,8.256)--(7.232,8.464)--(7.336,8.672)--(7.44,8.88)--(7.544,9.088)--(7.648,9.296)--(7.752,9.504)--(7.856,9.712)--(7.96,9.92)--(8.064,10.128)--(8.168,10.336)--(8.272,10.544)--(8.376,10.752)--(8.48,10.96)--(8.584,11.168)--(8.688,11.376)--(8.792,11.584)--(8.896,11.792)--(9,12);
\draw[green!55!black,line width=1.1pt,line join=round] (-4,1)--(-3.896,1.052)--(-3.792,1.104)--(-3.688,1.156)--(-3.584,1.208)--(-3.48,1.26)--(-3.376,1.312)--(-3.272,1.364)--(-3.168,1.416)--(-3.064,1.468)--(-2.96,1.52)--(-2.856,1.572)--(-2.752,1.624)--(-2.648,1.676)--(-2.544,1.728)--(-2.44,1.78)--(-2.336,1.832)--(-2.232,1.884)--(-2.128,1.936)--(-2.024,1.988)--(-1.92,2.04)--(-1.816,2.092)--(-1.712,2.144)--(-1.608,2.196)--(-1.504,2.248)--(-1.4,2.3)--(-1.296,2.352)--(-1.192,2.404)--(-1.088,2.456)--(-0.984,2.508)--(-0.88,2.56)--(-0.776,2.612)--(-0.672,2.664)--(-0.568,2.716)--(-0.464,2.768)--(-0.36,2.82)--(-0.256,2.872)--(-0.152,2.924)--(-0.048,2.976)--(0.056,3.028)--(0.16,3.08)--(0.264,3.132)--(0.368,3.184)--(0.472,3.236)--(0.576,3.288)--(0.68,3.34)--(0.784,3.392)--(0.888,3.444)--(0.992,3.496)--(1.096,3.548)--(1.2,3.6)--(1.304,3.652)--(1.408,3.704)--(1.512,3.756)--(1.616,3.808)--(1.72,3.86)--(1.824,3.912)--(1.928,3.964)--(2.032,4.016)--(2.136,4.068)--(2.24,4.12)--(2.344,4.172)--(2.448,4.224)--(2.552,4.276)--(2.656,4.328)--(2.76,4.38)--(2.864,4.432)--(2.968,4.484)--(3.072,4.536)--(3.176,4.588)--(3.28,4.64)--(3.384,4.692)--(3.488,4.744)--(3.592,4.796)--(3.696,4.848)--(3.8,4.9)--(3.904,4.952)--(4.008,5.004)--(4.112,5.056)--(4.216,5.108)--(4.32,5.16)--(4.424,5.212)--(4.528,5.264)--(4.632,5.316)--(4.736,5.368)--(4.84,5.42)--(4.944,5.472)--(5.048,5.524)--(5.152,5.576)--(5.256,5.628)--(5.36,5.68)--(5.464,5.732)--(5.568,5.784)--(5.672,5.836)--(5.776,5.888)--(5.88,5.94)--(5.984,5.992)--(6.088,6.044)--(6.192,6.096)--(6.296,6.148)--(6.4,6.2)--(6.504,6.252)--(6.608,6.304)--(6.712,6.356)--(6.816,6.408)--(6.92,6.46)--(7.024,6.512)--(7.128,6.564)--(7.232,6.616)--(7.336,6.668)--(7.44,6.72)--(7.544,6.772)--(7.648,6.824)--(7.752,6.876)--(7.856,6.928)--(7.96,6.98)--(8.064,7.032)--(8.168,7.084)--(8.272,7.136)--(8.376,7.188)--(8.48,7.24)--(8.584,7.292)--(8.688,7.344)--(8.792,7.396)--(8.896,7.448)--(9,7.5);
\filldraw[fill=blue!70!black,draw=blue!70!black,line width=.8pt] (6,6) circle (1.5pt);
\end{scope}\end{tikzpicture}\end{center}


\ex{2} $x=y^2-4\Rightarrow y^2=x+4\Rightarrow\res{f^{-1}(x)=\sqrt{x+4}}$. \ $D(f)=[0,\,+\infty)$, $R(f)=[-4,\,+\infty)$; \ $D(f^{-1})=\res{[-4,\,+\infty)}$, $R(f^{-1})=\res{[0,\,+\infty)}$: s'intercanvien. Sense la condició, dos valors de $x$ tindrien la mateixa imatge ($f(2)=f(-2)=0$) i la inversa no seria una funció.

\ex{3} \begin{enumerate}[label=\alph*),leftmargin=*,itemsep=2pt,topsep=2pt]\item $F(25)=45+32=\res{77}$ °F\item $f=1{,}8c+32\Rightarrow\res{C(f)=\dfrac{f-32}{1{,}8}}$: passa de graus Fahrenheit a graus Celsius.\item $C(212)=\frac{180}{1{,}8}=\res{100}$ °C; \ $C(50)=\frac{18}{1{,}8}=\res{10}$ °C\end{enumerate}


% =====================================================
\seccio{9}{Estudi complet d'una funció a partir de la gràfica}

\exf{2} $D=\res{[-5,\,6]}$, $R=\res{[-2,\,4]}$. Talls: $OX$: $\res{(-5,\,0),\ (-1,\,0),\ (1,\,0),\ (5,\,0)}$; $OY$: $\res{(0,\,-1)}$. Discontínua en $x=0$ (\rest{salt finit}: per l'esquerra arriba a $-1=f(0)$ i per la dreta comença en $-2$). Creixent en $\res{(-5,\,-3)\cup(0,\,3)}$; decreixent en $\res{(-3,\,0)\cup(3,\,6)}$. Màxim relatiu $\res{(-3,\,2)}$ i màxim absolut $\res{(3,\,4)}$; mínim absolut $\res{(6,\,-2)}$. No és simètrica ni periòdica.

\ex{2} \begin{enumerate}[label=\alph*),leftmargin=*,itemsep=2pt,topsep=2pt]\item $D=\res{[0,\,5]}$ i $R=\res{[0,\,40]}$: l'eixida dura 5 h i arriba a 40 km de casa (màxim absolut $(3,\,40)$).\item Creixent en $\res{(0,\,1)\cup(2,\,3)}$, constant en $\res{(1,\,2)}$ (està parada) i decreixent en $\res{(3,\,5)}$ (torna).\item Fa 40 km en 2 h: \rest{20 km/h}.\item Al cap de \rest{mitja hora} (anada) i de \rest{4 hores i mitja} (tornada).\end{enumerate}

\ex{3} Com que és senar, cada punt $(a,\,b)$ té el simètric $(-a,\,-b)$: el mínim és $\res{(2,\,-3)}$, $f(-4)=-1$ i $f(0)=\res{0}$. Per exemple:

\begin{center}\begin{tikzpicture}[x=0.42cm,y=0.42cm,baseline={(current bounding box.center)},font=\scriptsize]
\draw[muted!25,line width=.3pt] (-5,-4)--(-5,4);
\draw[muted!25,line width=.3pt] (-4,-4)--(-4,4);
\draw[muted!25,line width=.3pt] (-3,-4)--(-3,4);
\draw[muted!25,line width=.3pt] (-2,-4)--(-2,4);
\draw[muted!25,line width=.3pt] (-1,-4)--(-1,4);
\draw[muted!25,line width=.3pt] (0,-4)--(0,4);
\draw[muted!25,line width=.3pt] (1,-4)--(1,4);
\draw[muted!25,line width=.3pt] (2,-4)--(2,4);
\draw[muted!25,line width=.3pt] (3,-4)--(3,4);
\draw[muted!25,line width=.3pt] (4,-4)--(4,4);
\draw[muted!25,line width=.3pt] (5,-4)--(5,4);
\draw[muted!25,line width=.3pt] (-5,-4)--(5,-4);
\draw[muted!25,line width=.3pt] (-5,-3)--(5,-3);
\draw[muted!25,line width=.3pt] (-5,-2)--(5,-2);
\draw[muted!25,line width=.3pt] (-5,-1)--(5,-1);
\draw[muted!25,line width=.3pt] (-5,0)--(5,0);
\draw[muted!25,line width=.3pt] (-5,1)--(5,1);
\draw[muted!25,line width=.3pt] (-5,2)--(5,2);
\draw[muted!25,line width=.3pt] (-5,3)--(5,3);
\draw[muted!25,line width=.3pt] (-5,4)--(5,4);
\draw[-latex,muted] (-5,0)--(5,0);
\draw[-latex,muted] (0,-4)--(0,4);
\node[below=-1pt,text=muted] at (-5,0) {$-5$};
\node[below=-1pt,text=muted] at (-4,0) {$-4$};
\node[below=-1pt,text=muted] at (-3,0) {$-3$};
\node[below=-1pt,text=muted] at (-2,0) {$-2$};
\node[below=-1pt,text=muted] at (-1,0) {$-1$};
\node[below=-1pt,text=muted] at (1,0) {$1$};
\node[below=-1pt,text=muted] at (2,0) {$2$};
\node[below=-1pt,text=muted] at (3,0) {$3$};
\node[below=-1pt,text=muted] at (4,0) {$4$};
\node[below=-1pt,text=muted] at (5,0) {$5$};
\node[left=-1pt,text=muted] at (0,-4) {$-4$};
\node[left=-1pt,text=muted] at (0,-3) {$-3$};
\node[left=-1pt,text=muted] at (0,-2) {$-2$};
\node[left=-1pt,text=muted] at (0,-1) {$-1$};
\node[left=-1pt,text=muted] at (0,1) {$1$};
\node[left=-1pt,text=muted] at (0,2) {$2$};
\node[left=-1pt,text=muted] at (0,3) {$3$};
\node[left=-1pt,text=muted] at (0,4) {$4$};
\begin{scope}\clip (-5,-4) rectangle (5,4);
\draw[accent,line width=1.1pt,line join=round] (-4,-1)--(-2,3);
\draw[accent,line width=1.1pt,line join=round] (-2,3)--(2,-3);
\draw[accent,line width=1.1pt,line join=round] (2,-3)--(4,1);
\filldraw[fill=accent,draw=accent,line width=.8pt] (-4,-1) circle (1.5pt);
\filldraw[fill=accent,draw=accent,line width=.8pt] (4,1) circle (1.5pt);
\filldraw[fill=blue!70!black,draw=blue!70!black,line width=.8pt] (-2,3) circle (1.5pt);
\filldraw[fill=blue!70!black,draw=blue!70!black,line width=.8pt] (2,-3) circle (1.5pt);
\end{scope}\end{tikzpicture}\end{center}


\ex{3} \begin{tasks}(1)
  \task $D=\res{\mathbb{R}}$; $OX$: $\res{(-2,\,0),\ (0,\,0),\ (2,\,0)}$, $OY$: $\res{(0,\,0)}$; \rest{parella}.
  \task $D=\res{\mathbb{R}-\{-1,\,1\}}$; talla els dos eixos en $\res{(0,\,0)}$; \rest{senar}.
  \task $D=\res{[-2,\,2]}$; $OX$: $\res{(-2,\,0),\ (2,\,0)}$, $OY$: $\res{(0,\,2)}$; \rest{parella}.
\end{tasks}


\end{document}
